\subsection{\texorpdfstring{Homomorphisms \(\mathbf{B} \otimes_{\mathbf{A}} \operatorname{Tor}^{\mathbf{A}}(\mathbf{E}, \mathbf{F}) \to \operatorname{Tor}^{\mathbf{B}}(\mathbf{E} \otimes_{\mathbf{A}} \mathbf{B}, \mathbf{B} \otimes_{\mathbf{A}} \mathbf{F})\) and \(\mathbf{B} \otimes_{\mathbf{A}} \operatorname{Ext}_{\mathbf{A}}(\mathbf{E}, \mathbf{F}) \to \operatorname{Ext}_{\mathbf{B}}(\mathbf{B} \otimes_{\mathbf{A}} \mathbf{E}, \mathbf{B} \otimes_{\mathbf{A}} \mathbf{F})\)}{Extension-of-scalars homomorphisms for Tor and Ext}}

In this section, we assume that A is commutative; we are given a ring homomorphism \(\rho: A \to B\) such that \(\rho(A)\) be contained in the center of B, and let E and F be two A-modules. We have a canonical isomorphism of complexes of A-modules

\[
u: \mathrm{B} \otimes_ {\mathrm{A}} (\mathrm{L} (\mathrm{E}) \otimes_ {\mathrm{A}} \mathrm{L} (\mathrm{F})) \rightarrow (\mathrm{L} (\mathrm{E}) \otimes_ {\mathrm{A}} \mathrm{B}) \otimes_ {\mathrm{B}} (\mathrm{B} \otimes_ {\mathrm{A}} \mathrm{L} (\mathrm{F}));
\]

on the other hand, since L(E) ⊗\_A B and B ⊗\_If L(F) are free B-complexes, there is a canonical homomorphism of graded A-modules (X, p.~102)

\[
\begin{array}{r l} \psi^ {\prime} (\mathrm{L(E)} \otimes_ {\mathrm{A}} \mathrm{B}, \mathrm{B} \otimes_ {\mathrm{A}} \mathrm{L(F)}): & \mathrm{H} ((\mathrm{L(E)} \otimes_ {\mathrm{A}} \mathrm{B}) \otimes_ {\mathrm{B}} (\mathrm{B} \otimes_ {\mathrm{A}} \mathrm{L(F)})) \\ & \to \operatorname{Tor} ^ {\mathrm{B}} (\mathrm{E} \otimes_ {\mathrm{A}} \mathrm{B}, \mathrm{B} \otimes_ {\mathrm{A}} \mathrm{F}) \end{array}
\]

Finally, we have a homomorphism (X, p.~62)

\[
\gamma : \mathrm{B} \otimes_ {\mathbf {A}} \operatorname{Tor} ^ {\mathbf {A}} (\mathrm{E}, \mathrm{F}) \rightarrow \mathrm{H} (\mathrm{B} \otimes_ {\mathbf {A}} \mathrm{L} (\mathrm{E}) \otimes_ {\mathbf {A}} \mathrm{L} (\mathrm{F})) .
\]

The canonical homomorphism of graded B-modules

\[
\mathrm{B} \otimes_ {\mathrm{A}} \operatorname{Tor} ^ {\mathrm{A}} (\mathrm{E}, \mathrm{F}) \rightarrow \operatorname{Tor} ^ {\mathrm{B}} (\mathrm{E} \otimes_ {\mathrm{A}} \mathrm{B}, \mathrm{B} \otimes_ {\mathrm{A}} \mathrm{F})\tag{13}
\]

is defined as the composite \(\psi^{\prime}(\mathrm{L}(\mathrm{E})\otimes_{\mathrm{A}}\mathrm{B},\mathrm{B}\otimes_{\mathrm{A}}\mathrm{L}(\mathrm{F}))\circ \mathrm{H}(u)\circ \gamma .\)

PROPOSITION 9. --- If B is flat over A, then the homomorphism (13) is bijective.

Indeed, \(\psi^{\prime}(\mathbf{L}(\mathbf{E})\otimes_{\mathbf{A}}\mathbf{B},\mathbf{B}\otimes_{\mathbf{A}}\mathbf{L}(\mathbf{F}))\) is bijective (X, p.~102, prop. 1) and \(\gamma\) is bijective (X, p.~66, cor. 2).

Suppose B is flat over A. Substituting in the homomorphism (12) E for Q, B for N, and \(\mathbf{B} \otimes_{\mathbf{A}} \mathbf{F}\) for M, we obtain a homomorphism

\[
\operatorname{Ext} _ {\mathbf {A}} \left(\mathrm{E}, \mathrm{B} \otimes_ {\mathbf {A}} \mathrm{F}\right)\rightarrow \operatorname{Ext} _ {\mathbf {B}} \left(\mathrm{B} \otimes_ {\mathbf {A}} \mathrm{E}, \mathrm{B} \otimes_ {\mathbf {A}} \mathrm{F}\right)
\]

which is bijective by Prop. 8. Substituting E for M, F for N, A for B, B for Q in the homomorphism (10), and interchanging the factors of the tensor products, we obtain a homomorphism of B-modules

\[
\mathrm{B} \otimes_ {\mathrm{A}} \operatorname{Ext} _ {\mathrm{A}} (\mathrm{E}, \mathrm{F}) \rightarrow \operatorname{Ext} _ {\mathrm{A}} (\mathrm{E}, \mathrm{B} \otimes_ {\mathrm{A}} \mathrm{F}),
\]

whence by composition a homomorphism, called canonical

\[
\mathbf {B} \otimes_ {\mathbf {A}} \operatorname{Ext} _ {\mathbf {A}} (\mathbf {E}, \mathbf {F}) \rightarrow \operatorname{Ext} _ {\mathbf {B}} (\mathbf {B} \otimes_ {\mathbf {A}} \mathbf {E}, \mathbf {B} \otimes_ {\mathbf {A}} \mathbf {F}).\tag{14}
\]

PROPOSITION 10. --- The homomorphism (14) is bijective in the following cases:

\begin{enumerate}
\def\labelenumi{\alph{enumi})}
\item
  B is a finitely generated projective A-module;
\item
  B is a flat A-module, A is Noetherian, and E is a finitely generated A-module. This follows from prop. 7 (X, p.~108).
\end{enumerate}

